% sections/02-common-conventions.tex

\section{Common API Conventions}

\subsection{Base URL}

During local development, the backend is available at:

\begin{minted}{text}
http://localhost:3000
\end{minted}

All endpoint paths in this document are written relative to this base URL.

For example:

\begin{minted}{http}
GET /products
\end{minted}

means:

\begin{minted}{text}
http://localhost:3000/products
\end{minted}

\subsection{Content Type}

All request bodies should be sent as JSON.

\begin{minted}{http}
Content-Type: application/json
\end{minted}

Unless stated otherwise, endpoints that accept a request body expect JSON.

\subsection{Authentication}

Authenticated endpoints require a Bearer token:

\begin{minted}{http}
Authorization: Bearer <accessToken>
\end{minted}

The API uses three practical access levels:

\begin{itemize}
    \item \textbf{Public} --- no token required,
    \item \textbf{User} --- authenticated regular user,
    \item \textbf{Admin} --- authenticated administrator.
\end{itemize}

Some endpoints are available to both users and administrators, but apply ownership rules. For example, a regular user can access only their own cart, orders, payments, and reviews, while an administrator may have broader access for selected operations.

\subsection{Identifiers}

All persistent entity identifiers are MongoDB ObjectId strings.

Example:

\begin{minted}{text}
665000000000000000000001
\end{minted}

Path parameters named \texttt{id}, \texttt{productId}, \texttt{categoryId}, \texttt{orderId}, \texttt{paymentId}, or \texttt{reviewId} should be valid ObjectId strings.

Invalid ObjectId values are rejected before reaching the database layer.

\subsection{Pagination}

Paginated endpoints use the following query parameters:

\begin{itemize}
    \item \texttt{page} --- optional positive integer, default: \texttt{1},
    \item \texttt{limit} --- optional positive integer, default: \texttt{20}, maximum: \texttt{100}.
\end{itemize}

Example:

\begin{minted}{http}
GET /products?page=1&limit=20
\end{minted}

Paginated responses use the following general shape:

\begin{minted}[frame=single]{json}
{
"items": [],
"page": 1,
"limit": 20,
"total": 0,
"totalPages": 0
}
\end{minted}

\subsection{Money Representation}

Prices and monetary amounts are represented as integer-like numeric values in minor currency units.

For example:

\begin{minted}[frame=single]{json}
{
"amount": 29999,
"currency": "PLN"
}
\end{minted}

means:

\begin{minted}{text}
299.99 PLN
\end{minted}

The currently supported currencies are:

\begin{itemize}
    \item \texttt{PLN},
    \item \texttt{EUR},
    \item \texttt{USD}.
\end{itemize}

\subsection{Common Status and Type Values}

\subsubsection{User roles}

\begin{itemize}
    \item \texttt{user},
    \item \texttt{admin}.
\end{itemize}

\subsubsection{User statuses}

\begin{itemize}
    \item \texttt{active},
    \item \texttt{deleted}.
\end{itemize}

\subsubsection{Category statuses}

\begin{itemize}
    \item \texttt{active},
    \item \texttt{archived}.
\end{itemize}

\subsubsection{Product statuses}

\begin{itemize}
    \item \texttt{draft},
    \item \texttt{active},
    \item \texttt{archived}.
\end{itemize}

\subsubsection{Category attribute types}

\begin{itemize}
    \item \texttt{string},
    \item \texttt{number},
    \item \texttt{boolean},
    \item \texttt{string\_array}.
\end{itemize}

These types are used by category attribute definitions and determine how the frontend should render product attribute input fields.

\begin{table}[H]
    \centering
    \begin{tabular}{|l|l|}
        \hline
        \textbf{Attribute type}                               & \textbf{Suggested frontend input} \\
        \hline
        \texttt{string} with \texttt{allowedValues}           & select/dropdown                   \\
        \texttt{string} without \texttt{allowedValues}        & text input                        \\
        \texttt{number}                                       & number input                      \\
        \texttt{boolean}                                      & checkbox/toggle                   \\
        \texttt{string\_array} with \texttt{allowedValues}    & multi-select                      \\
        \texttt{string\_array} without \texttt{allowedValues} & tag input                         \\
        \hline
    \end{tabular}
    \caption{Frontend interpretation of category attribute types}
\end{table}

\subsubsection{Order statuses}

\begin{itemize}
    \item \texttt{pending\_payment},
    \item \texttt{paid},
    \item \texttt{payment\_failed},
    \item \texttt{cancelled},
    \item \texttt{shipped},
    \item \texttt{completed}.
\end{itemize}

\subsubsection{Payment statuses}

\begin{itemize}
    \item \texttt{pending},
    \item \texttt{paid},
    \item \texttt{failed}.
\end{itemize}

\subsubsection{Review statuses}

\begin{itemize}
    \item \texttt{active},
    \item \texttt{deleted}.
\end{itemize}

\subsection{Common Error Response}

Validation and business rule errors usually return a JSON response with a message, HTTP error name, and status code.

Example:

\begin{minted}[frame=single]{json}
{
"message": "Product not found",
"error": "Not Found",
"statusCode": 404
}
\end{minted}

For DTO validation errors, the \texttt{message} field may be an array:

\begin{minted}[frame=single]{json}
{
"message": [
"email must be an email",
"password must be longer than or equal to 6 characters"
],
"error": "Bad Request",
"statusCode": 400
}
\end{minted}

\subsection{Database Validation Errors}

Some write operations may fail MongoDB schema validation. In that case, the API returns a response with additional \texttt{validationErrors} details.

Example:

\begin{minted}[frame=single]{json}
{
"message": "Category failed database validation",
"validationErrors": [
"attributeDefinitions failed validation"
],
"error": "Bad Request",
"statusCode": 400
}
\end{minted}

These errors usually indicate a mismatch between the service layer and the MongoDB collection validator. They are useful during development because they show which database-level rule rejected the document.

\subsection{Soft Deletion and Archiving}

The API generally avoids hard deletion for business entities.

\begin{itemize}
    \item Users are soft-deleted by changing their status to \texttt{deleted}.
    \item Reviews are soft-deleted by changing their status to \texttt{deleted}.
    \item Products and categories are archived by changing their status to \texttt{archived}.
\end{itemize}

Public endpoints usually return only active records. Admin endpoints may expose inactive or archived records where needed.

\subsection{Timestamps}

Date-time fields are returned as ISO strings.

Common timestamp fields include:

\begin{itemize}
    \item \texttt{createdAt},
    \item \texttt{updatedAt},
    \item \texttt{deletedAt},
    \item \texttt{archivedAt},
    \item \texttt{paidAt},
    \item \texttt{failedAt},
    \item \texttt{cancelledAt}.
\end{itemize}

Optional timestamp fields are returned only when relevant.

\subsection{Endpoint Documentation Format}

Each endpoint is described using the following structure:

\begin{itemize}
    \item \textbf{Access} --- public, authenticated user, administrator, or mixed ownership-based access.
    \item \textbf{Description} --- what the endpoint does.
    \item \textbf{Query parameters} --- optional URL parameters, if supported.
    \item \textbf{Path parameters} --- required route parameters.
    \item \textbf{Request body} --- JSON body accepted by POST and PATCH endpoints.
    \item \textbf{Important rules} --- business rules or validation rules worth knowing.
\end{itemize}
